\renewcommand{\refname}{Список литературы}
\begin{thebibliography}{9}
\bibitem{zhang2006} Zhang W., Govorov A.~O., Bryant G.~W.
Semiconductor-Metal Nanoparticle Molecules: Hybrid Excitons and the
Nonlinear Fano Effect // Physical Review Letters. 2006. Vol.~97. 146804.
\url{https://doi.org/10.1103/PhysRevLett.97.146804}.

\bibitem{shah2013} Shah R.~A., Scherer N.~F., Pelton M., Gray S.~K.
Ultrafast reversal of a Fano resonance in a plasmon-exciton system
// Physical Review B. 2013. Vol.~88. 075411.
\url{https://doi.org/10.1103/PhysRevB.88.075411}.
Открытый текст: \url{https://arxiv.org/abs/1304.4538}.

\bibitem{johnson1972} Johnson P.~B., Christy R.~W.
Optical Constants of the Noble Metals // Physical Review B.
1972. Vol.~6. P.~4370--4379.
\url{https://doi.org/10.1103/PhysRevB.6.4370}.

\bibitem{meier1983} Meier M., Wokaun A.
Enhanced fields on large metal particles: dynamic depolarization
// Optics Letters. 1983. Vol.~8. P.~581--583.
\url{https://doi.org/10.1364/OL.8.000581}.

\bibitem{novo2006} Novo C., Gomez D., Perez-Juste J., Zhang Z.,
Petrova H., Reismann M., Mulvaney P., Hartland G.~V.
Contributions from radiation damping and surface scattering to the
linewidth of the longitudinal plasmon band of gold nanorods:
a single particle study // Physical Chemistry Chemical Physics.
2006. Vol.~8. P.~3540--3546.
\url{https://doi.org/10.1039/B604856K}.
\end{thebibliography}
